\documentclass[a4paper,12pt]{article}
\usepackage[utf8]{inputenc}
\usepackage[T1]{fontenc}
\usepackage[italian]{babel}
\usepackage{amsmath,amssymb}
\usepackage{booktabs}
\usepackage{geometry}
\geometry{margin=2.5cm}

\begin{document}

\section*{Formulazione QUBO del problema Knapsack}

Siano:
\[
x_i \in \{0,1\}, \qquad i=1,\dots,n,
\]
dove \(x_i = 1\) indica che l'oggetto \(i\) viene selezionato, mentre \(x_i = 0\) indica il contrario.

Indichiamo con:
\[
p_i \text{ il profitto dell'oggetto } i, \qquad
w_i \text{ il peso dell'oggetto } i, \qquad
C \text{ la capacità dello zaino.}
\]

\paragraph{}
Dove \(\lambda\) è stato fissato come rapporto tra la somma dei profitti degli oggetti e la capacità dello zaino, ossia:
\[
\lambda = \frac{\sum_{i=1}^{n} p_i}{C}.
\]
oppure
\[
\lambda = \frac{\sum_{i=1}^{n} p_i}{3}
\]

\paragraph{}
L'obiettivo è massimizzare il profitto totale rispettando il vincolo sulla capacità:
\[
\max \sum_{i=1}^{n} p_i x_i
\qquad \text{s.t.} \qquad
\sum_{i=1}^{n} w_i x_i = C.
\]

Per ottenere una formulazione QUBO, il problema viene riscritto come minimizzazione introducendo un termine di penalità quadratico:
\begin{equation}
Q(x) = -\sum_{i=1}^{n} p_i x_i + \lambda \left( \sum_{i=1}^{n} w_i x_i - C \right)^2 .
\end{equation}

Sviluppando il quadrato si ottiene:
\begin{equation}
Q(x) =
-\sum_{i=1}^{n} p_i x_i
+ \lambda \left[
\left(\sum_{i=1}^{n} w_i x_i\right)^2
- 2C \sum_{i=1}^{n} w_i x_i
+ C^2
\right].
\end{equation}

Espandendo il termine quadratico:
\begin{equation}
\left(\sum_{i=1}^{n} w_i x_i\right)^2
=
\sum_{i=1}^{n} w_i^2 x_i^2
+
2\sum_{i<j} w_i w_j x_i x_j.
\end{equation}

Poiché le variabili sono binarie, vale la proprietà
\[
x_i^2 = x_i.
\]

Pertanto:
\begin{equation}
Q(x) =
-\sum_{i=1}^{n} p_i x_i
+ \lambda \sum_{i=1}^{n} w_i^2 x_i
+ 2\lambda \sum_{i<j} w_i w_j x_i x_j
- 2\lambda C \sum_{i=1}^{n} w_i x_i
+ \lambda C^2.
\end{equation}

Trascurando il termine costante \(\lambda C^2\), che non influenza la posizione del minimo, si ottiene:
\begin{equation}
Q(x) \sim
-\sum_{i=1}^{n} p_i x_i
+ \lambda \sum_{i=1}^{n} w_i^2 x_i
- 2\lambda C \sum_{i=1}^{n} w_i x_i
+ 2\lambda \sum_{i<j} w_i w_j x_i x_j.
\end{equation}

Raccogliendo i termini lineari e quadratici:
\begin{equation}
Q(x) \sim
\sum_{i=1}^{n} x_i \left(-p_i + \lambda w_i^2 - 2\lambda C w_i \right)
+
\sum_{i<j} x_i x_j \left(2\lambda w_i w_j\right).
\end{equation}

Da questa espressione si ricava la matrice QUBO \(Q = [Q_{ij}]\), definita da:
\begin{equation}
Q_{ij} =
\begin{cases}
-p_i + \lambda w_i^2 - 2\lambda C w_i, & i=j, \\[6pt]
\lambda w_i w_j, & i\neq j.
\end{cases}
\end{equation}

La formulazione finale del problema è quindi:
\begin{equation}
\min \; x^T Q x.
\end{equation}

\paragraph{}
Si osservi che questa formulazione utilizza un termine di penalità quadratico che favorisce soluzioni 
il cui peso totale sia vicino alla capacità \(C\). Di conseguenza, essa rappresenta una penalizzazione 
\emph{soft} del vincolo \(\sum_i w_i x_i \leq C\), e non una sua imposizione esatta.

\subsection*{Istanza sperimentale \texttt{ksp\_1.txt} utilizzando QALS}

Si consideri la seguente istanza del problema knapsack:

\[
n = 10, \qquad C = 101
\]

con oggetti descritti dalle coppie \((w_i, p_i)\):
\[
(42,81),\,
(42,42),\,
(68,56),\,
(68,25),\,
(77,14),\,
(57,63),\,
(17,75),\,
(19,41),\,
(94,19),\,
(34,12).
\]

L'istanza è stata valutata eseguendo QALS per \(1000\) run. I risultati ottenuti sono i seguenti:

\begin{table}[h!]
\centering
\begin{tabular}{lc}
\toprule
Numero totale di run & 1000 \\
Best Weight & 101 \\
Best Profit & 198 \\
Avg Weight GAP & -74.4 \\
Avg Profit GAP & 13.7\% \\
\bottomrule
\end{tabular}
\caption{Risultati con $\lambda = \displaystyle \frac{\sum_{i=1}^{n} p_i}{3}$}
\end{table}

\begin{table}[h!]
\centering
\begin{tabular}{lc}
\toprule
Numero totale di run & 1000 \\
Best Weight & 101 \\
Best Profit & 198 \\
Avg Weight GAP & -76.5 \\
Avg Profit GAP & 13.8\% \\
\bottomrule
\end{tabular}
\caption{Risultati con $\lambda = \displaystyle \frac{\sum_{i=1}^{n} p_i}{C}$}
\end{table}

\subsection*{Discussione dei risultati su \texttt{ksp\_1.txt}}

In \(1000\) esecuzioni sull'istanza \texttt{ksp\_1.txt}, QALS ha prodotto circa 
gli stessi risulatati utlizzando due $\lambda$ diversi.

Il profit gap medio pari al \(13.75\%\) mostra che 
il profitto ottenuto risulta circa 1/10 rispetto a quello ottimo. Parallelamente, 
il weight gap medio di \(-75.5\) unità indica che le soluzioni trovate tendono a 
superare la capacità dello zaino.

\subsection*{Istanza sperimentale \texttt{ksp\_2.txt} utilizzando QALS}

Si consideri la seguente istanza del problema knapsack:
\[
n = 50, \qquad C = 1001
\]
con oggetti descritti dalle coppie \((w_i, p_i)\) riportate nel file \texttt{ksp\_2.txt}.

L'istanza è stata valutata eseguendo QALS per \(1000\) run.
\begin{table}[h!]
\centering
\begin{tabular}{lc}
\toprule
Numero totale di run & 1000 \\
Best Weight & 978 \\
Best Profit & 7740 \\
Avg Weight GAP & -8604.6 \\
Avg Profit GAP & -74.1\% \\
\bottomrule
\end{tabular}
\caption{Risultati con $\lambda = \displaystyle \frac{\sum_{i=1}^{n} p_i}{3}$}
\end{table}

\begin{table}[h!]
\centering
\begin{tabular}{lc}
\toprule
Numero totale di run & 1000 \\
Best Weight & 978 \\
Best Profit & 7740 \\
Avg Weight GAP & -8646.7 \\
Avg Profit GAP & -75\% \\
\bottomrule
\end{tabular}
\caption{Risultati con $\lambda = \displaystyle \frac{\sum_{i=1}^{n} p_i}{C}$}
\end{table}

\subsection*{Discussione dei risultati su \texttt{ksp\_2.txt}}
In \(1000\) esecuzioni sull'istanza \texttt{ksp\_1.txt}, QALS ha prodotto circa 
gli stessi risulatati utlizzando due $\lambda$ diversi.

Il profit gap medio pari al \(-74.5\%\) mostra che 
il profitto ottenuto risulta più della metà di quello ottimo. Parallelamente, 
il weight gap medio di \(-8626.65\) unità indica che le soluzioni trovate tendono a 
superare di molto la capacità dello zaino.

\subsection*{Istanza sperimentale \texttt{ksp\_1.txt} NON utilizzando QALS}
\begin{table}[h!]
\centering
\begin{tabular}{lc}
\toprule
Numero totale di run & 1000 \\
Best Weight & 101 \\
Best Profit & 198 \\
Avg Weight GAP & -0.6 \\
Avg Profit GAP & 42.2\% \\
\bottomrule
\end{tabular}
\caption{Risultati con $\lambda = \displaystyle \frac{\sum_{i=1}^{n} p_i}{3}$}
\end{table}

\begin{table}[h!]
\centering
\begin{tabular}{lc}
\toprule
Numero totale di run & 1000 \\
Best Weight & 101 \\
Best Profit & 198 \\
Avg Weight GAP & -0.8 \\
Avg Profit GAP & 41\% \\
\bottomrule
\end{tabular}
\caption{Risultati con $\lambda = \displaystyle \frac{\sum_{i=1}^{n} p_i}{C}$}
\end{table}

\subsection*{Discussione dei risultati su \texttt{ksp\_1.txt}}
Le conclusioni per qaunto riguarda lambda sono sempre le stesse: non cambia molto utlizzare
una versione piuttosto di un'altra.
Per quanto riguarda i pesi e i profitti si può notare come le soluzioni che si trovano 
tendono a definre un gap nel peso minore.

\subsection*{Istanza sperimentale \texttt{ksp\_2.txt} NON utilizzando QALS}

\begin{table}[h!]
\centering
\begin{tabular}{lc}
\toprule
Numero totale di run & 1000 \\
Best Weight & 978 \\
Best Profit & 7740 \\
Avg Weight GAP & -24.5 \\
Avg Profit GAP & 53.0\% \\
\bottomrule
\end{tabular}
\caption{Risultati con $\lambda = \displaystyle \frac{\sum_{i=1}^{n} p_i}{3}$}
\end{table}

\begin{table}[h!]
\centering
\begin{tabular}{lc}
\toprule
Numero totale di run & 1000 \\
Best Weight & 978 \\
Best Profit & 7740 \\
Avg Weight GAP & -25.1 \\
Avg Profit GAP & 53.5\% \\
\bottomrule
\end{tabular}
\caption{Risultati con $\lambda = \displaystyle \frac{\sum_{i=1}^{n} p_i}{C}$}
\end{table}

\subsection*{Discussione dei risultati su \texttt{ksp\_2.txt}}
Le conclusioni per qaunto riguarda lambda sono sempre le stesse: non cambia molto utlizzare
una versione piuttosto di un'altra.
Per quanto riguarda i pesi e i profitti si può notare come le soluzioni che si trovano 
tendono a definre un gap nel peso minore.
Si noti come aumentando di complessità il problema il \(Avg\ Profit\_ GAP\) risulti maggiore.

\end{document}